\documentclass[book.tex]{subfiles}
\begin{document}

\chapter{Lorentz Group Equivariant Networks}

\label{chap:lorentz_equivariant}
\noindent\rule{11cm}{0.4pt} \\
\textbf{Prerequisites:} \autoref{chap:mpnn} \\
\textbf{Difficulty Level:} ** \\
\noindent\rule{11cm}{0.4pt}
\vspace{5mm}

We have studied a number of group equivariant architectures. For applications in high energy physics, it becomes very useful to have a Lorentz group equivariant architecture \cite{bogatskiy2020lorentz}. This architecture leverages the theory of finite dimensional representations of the Lorentz group. The resulting equivariant architecture is able to perform classification tasks from high energy physics with fewer datapoints than comparable baselines.

%"""
%We present a neural network architecture that is fully equivariant with respect to transformations under the Lorentz group, a fundamental symmetry of space and time in physics. The architecture is based on the theory of the finite-dimensional representations of the Lorentz group and the equivariant nonlinearity involves the tensor product. For classification tasks in particle physics, we demonstrate that such an equivariant architecture leads to drastically simpler models that have relatively few learnable parameters and are much more physically interpretable than leading approaches that use CNNs and point cloud approaches. The competitive performance of the network is demonstrated on a public classification dataset [27] for tagging top quark decays given energy-momenta of jet constituents produced in proton-proton collisions.
%""" 

\section{A Brief Review of Scattering}

The electron-muon scattering element is given by

\begin{align}
    \mathcal{M}^2 &= \frac{g_c^4[p_1\cdot p_3 +m_e^2][p_2\cdot p_4 + m_{\mu}^2]}{((p_1-p_2)^2 -m_{\gamma}^2)^2}
\end{align}

\section{A Review of the Lorentz Group}

The finite dimensional irreducible representations of the Lorentz group are given by tensor products of representations of $SU(2)$.
\begin{align}
    T^{(k,n)} &= T^{(k, 0)} \oplus T^{(0, n)}
\end{align}

The proper orthochronous group $SO(1,3)^+$ is isomorphic to $PSL(2, \mathbb{C})$ %(\textcolor{red}{TODO: Indroduce the projective special linear group}).

The Clebsch-Gordon product in this case is given by
\begin{align}
    H: \bigoplus_{k, n} \tilde{T}^{(k,n)} \to T^{(k_1, n_1)} \otimes T^{(k_2, n_2)}
\end{align}


\section{A Review of Equivariance}

The nonlinearity in our networks is given by the Clebsch-Gordon coefficients. Suppose that $R_{l_1}$ and $R_{l_2}$ be irreducible representations of $SU(2)$. We can write

\begin{align}
    R_{l_1} \otimes R_{l_2} &\sim \oplus_{l} \tilde{R_l}
\end{align}

The Clebsch-Gordon isomorphism is given by this map

\begin{align}
    B: \bigoplus_{l=|l_1-l_2|}^{l_1+l_2} \to R_{l_1} \times R_{l_2}
\end{align}

For $SU(2)$ define the canonical basis $e_{l,m}$. The product space $R_{l_1}\otimes R_{l_2}$ has a canonical basis given by $e_{l_1,m_1} \otimes e_{l_2, m_2}$. The Clebsch-Gordon coefficients are the components of the map
\begin{align}
    B: \tilde{e}_{l,m} \mapsto \sum_{m_1, m_2} B^{l_1,m_1;l_2,m_2}_{l,m} e_{l_1,m_1} \otimes e_{l_2,m_2}
\end{align}

The inverse transformation is given by 

\begin{align}
    B^{-1}: R_{l_1} \otimes R_{l_2} \to \sum_{l} \tilde{R}_l
\end{align}
\begin{align}
    e_{l_1, m_1} \otimes e_{l_2, m_2} &= \sum (B^{-1})^{l,m}_{l_1,m_1;l_2,m_2} \tilde{e}_{l,m}
\end{align}
where $B^{-1} = B^T$. 

%\textcolor{red}{TODO: I'm still quite confused by how the Clebsch-Gordon transformations are implemented.}

\section{Lorentz Group Network}

The inputs are the 4 momenta of $N_{obj}$ particles from some collision event along with any associated scalar labels. The first layer is a linear layer that produces output features that reside in an irreducible representation for the Lorentz group. %(\textcolor{red}{TODO: Clear up the explanation here!})

%\textcolor{red}{TODO: What is a CG layer?} 
We let $\textrm{CG}$ denote the Clebsch-Gordon operator. The update rule is given by %(\textcolor{red}{TODO: Explain this formula})

\begin{align}
\mathcal{F}_i^{(p+1)} &= L_{CG}(\mathcal{F}^p)_i \\
&= W \cdot \left ( \mathcal{F}^p_i \oplus \textrm{CG}[\mathcal{F}^p_i]^{\otimes 2} \oplus \textrm{CG}[ \sum_j f(p_{ij}^2)p_{ij} \otimes \mathcal{F}^p_j] \right ) 
\end{align}

Here we have $p_{ij} = p_i - p_j$. The output is given by

\begin{align}
    w_{out} &= W_{out} \left ( \sum_i \mathcal{F}_i^{(N_{CG})} \right )
\end{align}

\section{Jet Tagging}

High energy quarks emit cascading gluon emissions. The machine learning challenge we face is classifying which quark emissions are interesting and worthy of further study.

\end{document}
